\documentclass[11pt,a4paper]{article}
\usepackage[utf8]{inputenc}
\usepackage[T1]{fontenc}
\usepackage[brazilian]{babel}
\usepackage[margin=2.3cm]{geometry}
\usepackage{graphicx}
\usepackage{booktabs}
\usepackage{array}
\usepackage{xcolor}
\usepackage{hyperref}
\usepackage{caption}
\usepackage{amsmath}
\usepackage{longtable}
\usepackage{float}
\usepackage{enumitem}

\hypersetup{
  colorlinks=true,
  linkcolor=blue!50!black,
  urlcolor=blue!50!black,
  citecolor=blue!50!black
}

\title{Detecção de Anomalias em Cabiúnas via AutoML\\
\large Experimento EXP6 --- Vibração de mancais como preditor\\
de alarmes de TC382\_03\_A e T5\_AVG\_A}
\author{Pipeline \texttt{cnn1d\_ae} / AutoML --- projeto \texttt{TesteMLCab}}
\date{Agosto de 2026}

\begin{document}
\maketitle

\begin{abstract}
Este relatório documenta o EXP6: a extensão do pipeline AutoML de detecção
de anomalias do Cabiúnas para incorporar sinais de \textbf{vibração de
mancais} (10 canais, 5 mancais $\times$ radial X/Y) como preditores dos
alarmes de \texttt{TC382\_03\_A} e \texttt{T5\_AVG\_A}, usando uma nova
fonte de dados com 2,5 anos de histórico (2024--2026). Documentamos passo a
passo a metodologia (máscaras, split temporal, grid search), um bug real
encontrado e corrigido durante o experimento (estado operacional não
excluía o próprio treino), e o resultado final: um \texttt{IsolationForest}
multivariado que detecta \textbf{65\% dos alarmes reais com antecedência
mediana de 12,4 horas}, com falso alerta de apenas \textbf{0,06\%}
(3 anomalias/dia) -- uma melhoria de duas ordens de grandeza em falso
positivo sobre o candidato ``vencedor'' automático do próprio grid search,
que sofre de um desequilíbrio na função de score.
\end{abstract}

\tableofcontents

% =====================================================================
\section{Objetivo}
% =====================================================================

\textbf{A ideia central deste experimento é usar os sinais de vibração dos
5 mancais da turbina como preditores antecipados de falha dos sensores de
temperatura \texttt{TC382\_03\_A} e \texttt{T5\_AVG\_A}.} A hipótese física
é que uma degradação mecânica no mancal (vibração anormal) pode preceder,
em horas, a manifestação térmica que hoje dispara o alarme -- ou seja, a
vibração funcionaria como um sinal de alerta precoce, antes que a
temperatura em si saia da faixa normal.

Concretamente, o experimento treina um único modelo multivariado que
recebe \texttt{TC382\_03\_A} + \texttt{T5\_AVG\_A} + os 10 canais de
vibração como entrada, mas é avaliado \emph{apenas} pela sua capacidade de
antecipar os alarmes reais desses dois sensores de temperatura (a vibração
entra só como informação de contexto/apoio -- ver \texttt{eval\_sensors} na
Seção~\ref{sec:metodologia}). Comparamos o resultado com o EXP5 (ver
\texttt{docs/analise\_automl\_exp5.md}), que usava só
\texttt{TC382\_03\_A} isolado, sem vibração.

% =====================================================================
\section{Dados e série temporal utilizada}
\label{sec:dados}
% =====================================================================

\subsection{Nova fonte}

Dataset ClearML \textbf{\texttt{Cabiunas full 2024-2026 30s}}
(\texttt{TesteMLCab}, ID \texttt{a97ba56ba14840fbb1125c2a82f883c9}):

\begin{itemize}
  \item \texttt{sensores\_full\_2024\_2026\_30s.csv}: \textbf{2024-01-01 a
        2026-04-30}, resolução de 30s ($\sim$2,45 milhões de linhas), 39
        colunas de sensor, completude de 99,5--100\%.
  \item \texttt{alarmes\_selecionados\_turbina\_a.csv}: 2022-01-04 a
        2026-04-18, 6.851 linhas (pares onset/clear \texttt{ACT}/\texttt{INACT}),
        47 tags, sem duplicação de linha (diferente do arquivo usado no
        EXP1--EXP5).
\end{itemize}

\subsection{Vibração disponível}

10 canais \texttt{TV\_35[1-5][X|Y]\_A}, confirmados pela descrição dos
próprios alarmes do arquivo (ex.: \texttt{TV\_351X\_A} = ``Vibração X
Mancal 1''):

\begin{center}
\begin{tabular}{@{}lll@{}}
\toprule
Mancal & Radial X & Radial Y \\
\midrule
1 & \texttt{TV\_351X\_A} & \texttt{TV\_351Y\_A} \\
2 & \texttt{TV\_352X\_A} & \texttt{TV\_352Y\_A} \\
3 & \texttt{TV\_353X\_A} & \texttt{TV\_353Y\_A} \\
4 & \texttt{TV\_354X\_A} & \texttt{TV\_354Y\_A} \\
5 & \texttt{TV\_355X\_A} & \texttt{TV\_355Y\_A} \\
\bottomrule
\end{tabular}
\end{center}

\subsection{Janela usada neste experimento}

Do período total disponível, o EXP6 usa \textbf{2024-07-01 em diante}
(\texttt{DATA\_START\_DATE}), dividido em:

\begin{itemize}
  \item \textbf{Treino/fit}: 2024-07-01 a 2025-06-30 (12 meses)
  \item \textbf{Avaliação out-of-sample (OOS)}: 2025-07-01 a 2026-04-30
        ($\sim$10 meses) -- alarmes e pontos que o modelo nunca viu
\end{itemize}

\begin{figure}[H]
\centering
\includegraphics[width=\textwidth]{00_overview_TC382_03_A.png}
\caption{Série bruta de \texttt{TC382\_03\_A}, jul/2024 a abr/2026. Sombra
cinza = \texttt{off\_longo} (parada prolongada); linha preta vertical =
início do período de avaliação OOS; linhas verdes tracejadas = alarmes
reais no período OOS; pontos vermelhos = anomalias detectadas pelo modelo
final (\texttt{iforest}, ver Seção~\ref{sec:resultado}).}
\end{figure}

\begin{figure}[H]
\centering
\includegraphics[width=\textwidth]{00_overview_T5_AVG_A.png}
\caption{Mesma visualização para \texttt{T5\_AVG\_A}.}
\end{figure}

% =====================================================================
\section{Metodologia -- passo a passo}
\label{sec:metodologia}
% =====================================================================

A pipeline AutoML (\texttt{src/cnn1d\_ae/automl\_pipeline.py}) opera
\textbf{ponto a ponto}: cada instante de tempo é uma amostra independente,
sem janelamento temporal (diferente do CNN-1D sequencial usado no resto do
projeto). Para o EXP6, o processamento de cada grupo de sensores segue
estes passos, nesta ordem:

\begin{enumerate}[leftmargin=*]

\item \textbf{Recorte temporal da série} (\texttt{DATA\_START\_DATE}):
  a série completa (jan/2024--abr/2026) é cortada para começar em
  01/07/2024, antes de qualquer outro processamento. Campo novo, criado
  para este experimento -- antes só existia corte de \emph{fim} de treino
  via \texttt{AUTOML\_OOS\_SPLIT\_DATE}.

\item \textbf{Normalização de schema do arquivo de alarmes}: o arquivo
  novo usa colunas diferentes do antigo (\texttt{"Tag Alarme"} em vez de
  \texttt{"Tag"}) e traz pares onset/clear (\texttt{Status} =
  \texttt{ACT/UNACK} ou \texttt{INACT/UNACK}) para o mesmo evento. Mantemos
  só as linhas \texttt{ACT} (onset) -- o clear não é algo que faça sentido
  ``prever com antecedência''.

\item \textbf{Construção do dataframe do grupo multivariado}: as colunas
  dos sensores do grupo (\texttt{TC382\_03\_A}, \texttt{T5\_AVG\_A} + 10
  canais \texttt{TV\_}) são unidas num único dataframe alinhado por
  timestamp, com interpolação de gaps curtos ($\leq$ 3 amostras) e
  \texttt{ffill}/\texttt{bfill} nas bordas.

\item \textbf{Features derivadas} (\texttt{ENABLE\_DERIVED\_FEATURES=true}):
  para cada um dos 12 sensores, são adicionadas 3 colunas -- média móvel
  (\texttt{roll\_med}), desvio padrão móvel (\texttt{roll\_std}) e a
  diferença ponto a ponto (\texttt{delta\_1}), janela de 12 amostras
  (6 minutos). Total: $12 \times 4 = 48$ features de entrada para o modelo.
  \textbf{Implementado neste experimento} -- antes só existia para grupos
  de um sensor só.

\item \textbf{Estado operacional} (\texttt{RUNNING\_A}): o sensor
  \texttt{NGP\_A}, usado como referência operacional em todos os
  experimentos anteriores, não existe no arquivo novo. Passamos a usar
  \texttt{RUNNING\_A} (flag 0/1), após confirmar que a limpeza padrão já
  existente no pipeline (\texttt{to\_numeric} + interpolação) resolve os
  valores sujos do historiador (1,49\% de strings tipo
  \texttt{\{'Name': 'No Data', ...\}}). O limiar de ``desligado''
  (\texttt{OFF\_ABS\_THRESHOLD}) foi ajustado de \texttt{5.0} (fazia
  sentido para \texttt{NGP\_A}, sensor de rotação contínuo) para
  \texttt{0.5} (limiar binário correto para um flag 0/1).

\item \textbf{Exclusão do conjunto de treino} -- um ponto é excluído do
  treino se \emph{qualquer} uma destas condições for verdadeira:
  \begin{itemize}
    \item está dentro de $\pm$24h (\texttt{EXCLUDE\_MINUTES\_AROUND\_ALARM=1440})
          de um alarme real de \texttt{TC382\_03\_A} ou \texttt{T5\_AVG\_A};
    \item está num gap longo (mais que \texttt{INTERPOLATE\_LIMIT} amostras
          consecutivas faltando);
    \item o estado operacional \textbf{não} é \texttt{"on"} -- ver achado
          na Seção~\ref{sec:bug}.
  \end{itemize}

\item \textbf{Clip de outliers}: cada coluna é limitada ao intervalo entre
  seus percentis 0,1\% e 99,9\% (\texttt{OUTLIER\_MODE=quantile}).

\item \textbf{Split temporal out-of-sample}: das amostras que sobraram após
  os filtros acima, só as anteriores a \texttt{AUTOML\_OOS\_SPLIT\_DATE}
  (2025-07-01) entram no ajuste do modelo e no cálculo do percentil de
  threshold.

\item \textbf{Normalização z-score}: média e desvio padrão calculados
  \emph{apenas} sobre o conjunto de treino (pré-split), aplicados a toda a
  série -- evita vazamento de informação do período de avaliação.

\item \textbf{Ajuste do modelo}: para cada um dos 3 tipos
  (\texttt{dense}/\texttt{ocsvm}/\texttt{iforest}), o modelo é treinado
  \emph{uma vez} sobre o conjunto normal de treino; o erro de reconstrução
  (ou score de anomalia) é calculado tanto no treino quanto em toda a
  série.

\item \textbf{Grid de threshold/debounce}: para cada modelo já treinado,
  varremos 7 percentis de threshold ($\{90, 93, 95, 97, 99, 99.5, 99.9\}$)
  $\times$ 6 valores de debounce ($\{1, 2, 3, 6, 12, 24\}$ pontos
  consecutivos exigidos) = 42 combinações por modelo, 126 no total.

\item \textbf{Avaliação por trial}: cada combinação gera um
  \texttt{hit\_rate} (fração de alarmes OOS com pelo menos um ponto anômalo
  dentro de $\pm$24h), um \texttt{normal\_alert\_rate} (fração de pontos
  "normais" -- em operação, longe de qualquer alarme -- marcados como
  anômalos) e um \texttt{composite\_score}:
  \[
  \text{balanced\_score} = \text{hit\_rate} - \text{fp\_penalty} \cdot \text{normal\_alert\_rate}^2
  \]
  com \texttt{fp\_penalty=2.0} neste experimento.

\item \textbf{Escolha do sensor-alvo da avaliação} (\texttt{eval\_sensors},
  campo novo criado para este experimento): os 10 canais de vibração
  entram no modelo só como \emph{feature} -- seus próprios alarmes não
  contam no \texttt{hit\_rate}/\texttt{composite\_score}, que fica restrito
  a \texttt{["TC382\_03\_A", "T5\_AVG\_A"]}.

\end{enumerate}

% =====================================================================
\section{Configuração do experimento}
% =====================================================================

\begin{table}[H]
\centering
\small
\begin{tabular}{@{}ll@{}}
\toprule
\textbf{Parâmetro} & \textbf{Valor} \\
\midrule
Grupo & \texttt{TC382\_T5\_vibracao\_mancais} \\
Sensores (features) & \texttt{TC382\_03\_A}, \texttt{T5\_AVG\_A} + 10 canais \texttt{TV\_} \\
\texttt{eval\_sensors} & \texttt{TC382\_03\_A}, \texttt{T5\_AVG\_A} \\
Janela de fit & 2024-07-01 a 2025-06-30 \\
Janela de avaliação OOS & 2025-07-01 a 2026-04-30 \\
\texttt{OPERATIONAL\_REF\_SENSOR} & \texttt{RUNNING\_A} (\texttt{OFF\_ABS\_THRESHOLD=0.5}) \\
\texttt{ENABLE\_DERIVED\_FEATURES} & \texttt{true} (janela 12 = 6 min) \\
Grid & 7 percentis $\times$ 6 debounces $\times$ 3 modelos = 126 trials \\
\texttt{AUTOML\_DENSE\_LAYERS} & [256, 128], dropout 0.1, 100 epochs, patience 15 \\
\texttt{AUTOML\_IFOREST\_N\_ESTIMATORS} & 200, \texttt{contamination}=0.05 \\
\texttt{AUTOML\_OCSVM\_MAX\_TRAIN\_SAMPLES} & 50.000 (subamostra -- RBF-SVM não escala p/ 380k+ pontos) \\
\texttt{AUTOML\_FP\_PENALTY} & 2.0 \\
\bottomrule
\end{tabular}
\caption{Configuração completa (\texttt{test\_grupo\_exp6\_vibracao.json}).}
\end{table}

% =====================================================================
\section{Achado pré-resultado: máscara operacional não cobria o treino}
\label{sec:bug}
% =====================================================================

Antes de aceitar qualquer resultado, verificamos se as máscaras cobriam de
fato o \emph{treino} e não só a \emph{avaliação}. Encontramos um problema
real: o estado operacional (\texttt{RUNNING\_A}) só zerava anomalias na
hora de \textbf{pontuar} -- nunca excluía os períodos desligados do próprio
\textbf{conjunto de treino}. O modelo estava sendo ajustado numa mistura de
dois regimes muito diferentes (turbina operando, $\sim$500--800°C, vs.
parada/fria, $\sim$25--40°C).

\textbf{Antes da correção} (task \texttt{bd84fa6c...}):
\texttt{hit\_rate}=97,5\% mas \texttt{normal\_alert\_rate}=10,3\% -- o
modelo aprendia sobretudo a detectar a própria transição liga/desliga, não
anomalias sutis em operação. A correção também resolveu de graça outro
problema: 99,9\% do platô de sensor ``morto'' em $-40{,}51^{\circ}$C de
\texttt{TC382\_03\_A} (5092 de 5098 pontos) coincide com
\texttt{RUNNING\_A==0}, então excluir ``não-on'' do treino remove esse
artefato sem precisar mexer no threshold de outlier separadamente.

\textbf{Correção}: o cálculo do estado operacional foi movido para
\emph{antes} da construção do conjunto de treino, e \texttt{state != "on"}
passou a ser mais um critério de exclusão (junto com proximidade de alarme
e gaps longos). Todos os resultados deste relatório já refletem essa
correção.

% =====================================================================
\section{Ranking de modelos e falsos positivos}
\label{sec:ranking}
% =====================================================================

\subsection{Vencedor pelo \texttt{composite\_score} não é o melhor candidato prático}

\begin{table}[H]
\centering
\resizebox{\textwidth}{!}{%
\begin{tabular}{@{}lccccc@{}}
\toprule
\textbf{Modelo} & \textbf{Percentil / debounce} & \textbf{hit\_rate} & \textbf{normal\_alert\_rate} & \textbf{anomalias/dia} & \textbf{composite\_score} \\
\midrule
ocsvm (vencedor oficial) & 97,0 / 1 & 97,5\% (39/40) & \textbf{26,15\%} & 730,6 & 0,946 \\
dense & 99,5 / 3 & 82,5\% (33/40) & 12,92\% & 376,4 & 0,931 \\
\textbf{iforest} & \textbf{99,9 / 6} & \textbf{75,0\% (30/40)} & \textbf{0,06\%} & \textbf{3,2} & 0,917 \\
\bottomrule
\end{tabular}%
}
\caption{Melhor trial de cada modelo, dos 126 do grid completo (task
\texttt{2c3c7f92e93441dc976acf7d0aeb906d}).}
\end{table}

O \texttt{composite\_score} do \texttt{ocsvm} vencedor (0,946) supera o do
\texttt{iforest} (0,917) mesmo com uma taxa de falso alerta \textbf{436
vezes maior} -- sinal de que \texttt{AUTOML\_FP\_PENALTY=2.0} não penaliza
falso positivo o suficiente para impedir essa escolha quando o
\texttt{hit\_rate} sobe muito. Confirmamos visualmente
(Figura~\ref{fig:tradeoff}): o \texttt{ocsvm} marca quase toda a faixa
operacional como anômala, sem discriminar nada.

\begin{figure}[H]
\centering
\includegraphics[width=0.85\textwidth]{03_tradeoff_scatter.png}
\caption{Trade-off detecção $\times$ falso alerta para os 126 trials do
grid. O \texttt{iforest} escolhido (estrela preta) está no canto
superior-esquerdo ideal (alta detecção, baixíssimo falso alerta); o
\texttt{ocsvm} vencedor pelo score (X vermelho) está bem à direita, com
falso alerta muito mais alto para um ganho marginal de detecção.}
\label{fig:tradeoff}
\end{figure}

\begin{figure}[H]
\centering
\includegraphics[width=0.95\textwidth]{02_fp_por_modelo.png}
\caption{Esquerda: falso alerta do melhor trial de cada modelo -- o
\texttt{iforest} é ordens de grandeza melhor. Direita: distribuição de
falso alerta considerando todos os 42 trials de cada modelo no grid --
mostra que o \texttt{iforest} tem mais configurações de baixo FP
disponíveis, não só o ponto ótimo isolado.}
\end{figure}

% =====================================================================
\section{Resultado final: Isolation Forest}
\label{sec:resultado}
% =====================================================================

\subsection{Configuração e métricas de calibração}

\begin{table}[H]
\centering
\begin{tabular}{@{}ll@{}}
\toprule
\textbf{Parâmetro} & \textbf{Valor} \\
\midrule
Modelo & \texttt{IsolationForest} (scikit-learn) \\
\texttt{contamination} & 0,05 \\
\texttt{n\_estimators} & 200 \\
\texttt{threshold\_percentile} & 99,9 \\
\texttt{threshold} (erro de reconstrução) & 0,652 \\
\texttt{debounce} & 6 pontos consecutivos \\
\midrule
\texttt{hit\_rate} (OOS) & 75,0\% (30/40) \\
\texttt{normal\_alert\_rate} & 0,06\% \\
\texttt{anomaly\_rate\_points\_per\_day} & 3,2 \\
\texttt{composite\_score} & 0,917 \\
\bottomrule
\end{tabular}
\caption{Candidato final. Reproduzido de forma independente numa segunda
run (\texttt{AUTOML\_MODELS=["iforest"]} apenas, task
\texttt{d79fbb7ce31d45298a0a01761eab8c50}) -- métricas e hiperparâmetros
idênticos aos do grid completo, confirmando reprodutibilidade.}
\end{table}

\subsection{Antecedência real: previsão vs. reação}

Como já formalizado na análise do EXP5, o \texttt{hit\_rate} usa uma janela
\emph{simétrica} de $\pm$24h -- não distingue detecção \emph{antes} do
alarme (previsão real) de detecção \emph{depois} (reação tardia).
Separando os 30 casos detectados:

\begin{table}[H]
\centering
\begin{tabular}{@{}lcc@{}}
\toprule
\textbf{Categoria} & \textbf{Qtd.} & \textbf{\% dos 40 alarmes} \\
\midrule
Detectados dentro de $\pm$24h & 30 & 75,0\% \\
\quad -- com \textbf{antecedência real} (antes do alarme) & \textbf{26} & \textbf{65,0\%} \\
\quad -- só depois/simultâneo (reação) & 4 & 10,0\% \\
Sem detecção & 10 & 25,0\% \\
\bottomrule
\end{tabular}
\caption{Categorias de detecção dos 40 alarmes OOS.}
\label{tab:categorias}
\end{table}

\textbf{Antecedência dos 26 casos preditivos:} mediana de \textbf{12,4h},
média de 12,4h, mínimo de 1,3h, máximo de 23,5h (sem bater no teto de 24h
da janela -- diferente do EXP5, onde quase metade dos casos preditivos
saturava o limite).

\subsection{Visão completa: os 40 alarmes, um por um}
\label{sec:paineis}

Em vez de mostrar só 1--2 exemplos escolhidos a dedo, as Figuras
\ref{fig:painel-pred}--\ref{fig:painel-miss} trazem \textbf{todos os 40
alarmes} do período de avaliação, agrupados pela categoria da
Tabela~\ref{tab:categorias} -- cada miniatura mostra $\pm$30h/6h ao redor
do alarme (linha verde tracejada), com os pontos vermelhos marcando onde o
modelo sinalizou anomalia. O sensor plotado em cada miniatura é o mesmo que
gerou aquele alarme específico (\texttt{TC382\_03\_A} ou
\texttt{T5\_AVG\_A}).

\begin{figure}[H]
\centering
\includegraphics[width=\textwidth]{04_painel_preditivos.png}
\caption{Os 26 casos com antecedência real, ordenados da maior para a
menor antecedência (23,5h $\to$ poucos minutos). Vários mostram um padrão
em ``escada'' -- quedas ou subidas sucessivas de temperatura, cada uma já
sinalizada como anomalia, horas antes do instante do alarme.}
\label{fig:painel-pred}
\end{figure}

\begin{figure}[H]
\centering
\includegraphics[width=0.7\textwidth]{05_painel_reativos.png}
\caption{Os 4 casos em que a detecção só veio depois (ou praticamente
junto) do alarme -- o modelo reage ao mesmo evento, sem antecipar.}
\label{fig:painel-react}
\end{figure}

\begin{figure}[H]
\centering
\includegraphics[width=\textwidth]{06_painel_perdidos.png}
\caption{Os 10 alarmes sem nenhuma detecção dentro de $\pm$24h. Note que a
maioria não mostra excursão visível na própria série de temperatura
plotada -- são alarmes cuja causa aparentemente não se manifesta (ao menos
não de forma óbvia) nem no sensor de temperatura alarmado, nem nos 30 min
antes/depois mostrados aqui.}
\label{fig:painel-miss}
\end{figure}

\subsection{Dois exemplos em destaque (ampliados)}

Para leitura mais detalhada, as Figuras~\ref{fig:zoom1}--\ref{fig:zoom2}
ampliam dois casos da Figura~\ref{fig:painel-pred}, mostrando
\texttt{TC382\_03\_A} e \texttt{T5\_AVG\_A} lado a lado (os dois sensores
se movem de forma muito correlacionada).

\begin{figure}[H]
\centering
\includegraphics[width=0.95\textwidth]{pred_01_20260225_2045.png}
\caption{Exemplo de detecção com antecedência de 23,5h antes do alarme de
25/02/2026. Padrão em ``escada'' -- quedas sucessivas de temperatura, cada
uma sinalizada como anomalia, culminando no alarme.}
\label{fig:zoom1}
\end{figure}

\begin{figure}[H]
\centering
\includegraphics[width=0.95\textwidth]{pred_03_20260408_2153.png}
\caption{Segundo exemplo de detecção com antecedência, em abril de 2026.}
\label{fig:zoom2}
\end{figure}

% =====================================================================
\section{Comparação com o EXP5}
% =====================================================================

\begin{table}[H]
\centering
\begin{tabular}{@{}lcc@{}}
\toprule
\textbf{Métrica} & \textbf{EXP5} (\texttt{TC382\_03\_A} só) & \textbf{EXP6} (+ vibração) \\
\midrule
\texttt{hit\_rate} & 65,0\% (26/40) & 75,0\% (30/40) \\
Antecedência real & 45,0\% (9/20) & 65,0\% (26/40) \\
\texttt{normal\_alert\_rate} & 2,43\% & \textbf{0,06\%} ($\sim$40x menor) \\
Mediana de antecedência & 13,4h & 12,4h \\
\bottomrule
\end{tabular}
\end{table}

\textbf{Ressalva:} as bases de alarme são diferentes (EXP5 usava o arquivo
antigo, com linhas duplicadas, deduplicadas manualmente para 20 alarmes
OOS; EXP6 usa o arquivo novo, onset-only, 40 alarmes OOS) -- não é uma
comparação perfeitamente equivalente. Ainda assim, a melhoria de $\sim$40x
em falso alerta é grande o suficiente para ser um sinal real de que
vibração + features derivadas + correção da máscara de treino ajudam.

% =====================================================================
\section{Limitações e próximos passos}
% =====================================================================

\begin{itemize}
  \item \textbf{\texttt{AUTOML\_FP\_PENALTY}} precisa subir (ex.: 2,0
        $\to$ 8--10) para o \texttt{composite\_score} parar de escolher
        candidatos de FP alto por padrão -- hoje só chegamos ao
        \texttt{iforest} inspecionando o ranking manualmente.
  \item \textbf{Variância de semente} do \texttt{iforest} vencedor ainda
        não foi checada neste experimento (foi no EXP5,
        \texttt{AUTOML\_SEED\_SWEEP\_N}) -- pendente antes de considerar
        definitivo.
  \item \textbf{Contribuição real da vibração} não isolada: o ganho pode
        vir da vibração em si, do \texttt{ENABLE\_DERIVED\_FEATURES}, ou da
        correção da máscara de treino (Seção~\ref{sec:bug}) -- os três
        mudaram juntos. Um grupo de controle (só TC382+T5, mesmos ajustes,
        sem vibração) isolaria o efeito.
  \item \textbf{Comparação com EXP5 não é perfeitamente equivalente}
        (bases de alarme diferentes -- ver Seção anterior).
  \item Amostra de 40 alarmes ainda é pequena; frações como 65\%/75\% têm
        incerteza estatística não desprezível.
\end{itemize}

\end{document}
